\chapter{Text}\label{chap:text}

This chapter covers the second block of the workshop:
how \hologo{LaTeX} turns the lines of the source file into paragraphs,
how to type quotation marks, dashes and other special characters,
lists, and the preamble with its packages.

\section{Spaces, line breaks and paragraphs}\label{sec:spaces}

\hologo{LaTeX} does not keep the line breaks of the source file.
It collects the words of a paragraph and breaks them into lines itself,
so that every line is filled evenly:

\slidecode{02-linebreaks.tex}
\sliderender{02-linebreaks.pdf}

Three lines in, one line out.
The rules behind this are few:

\textbf{A single line break} is a space, like the space between two words.
Several spaces in a row count as one.

\slidecode{02-newline.tex}
\sliderender{02-newline.pdf}

\textbf{An empty line} starts a new paragraph.
In the standard classes, a new paragraph is indented,
except for the first one after a heading.
\autoref{sec:parskip} shows how to get a gap between paragraphs instead.

\slidecode{02-paragraph.tex}
\sliderender{02-paragraph.pdf}

\textbf{A \texttt{\%} at the end of a line} comments out the line break, and with it the space.
This is mostly useful inside the definitions of your own commands (\autoref{sec:own-commands}),
where a stray space would end up in the output.

\slidecode{02-percent.tex}
\sliderender{02-percent.pdf}

\textbf{A tilde \texttt{\textasciitilde}} is a space at which the line is never broken.
Use it between a word and the number or name that belongs to it,
e.g.\ \code|Section~2|, \code|Figure~3|, \code|Dr.~Watson| or \code|10~K|.
With a normal space, the number can end up alone at the start of the next line:

\slidecode{02-space.tex}
\sliderender{02-space.pdf}

With the tilde, \enquote{Section} moves down to join it:

\slidecode{02-tilde.tex}
\sliderender{02-tilde.pdf}

\subsection{One sentence per line}

Since a single line break is just a space, how you break the lines in the source file is up to you.
A good habit is one sentence per line:

\slidecode{02-one-sentence.tex}

The PDF is the same, but the source is easier to work with.
Version control such as Git compares files line by line,
so it then shows the sentence that changed instead of the whole paragraph,
and an editor's \enquote{go to line} lands on the right sentence.
The workshop template and this handout are written this way.

\subsection{Forced line breaks}

\code|\\| ends a line at that point:

\slidecode{02-forced-break.tex}

This is almost never what you want in the text of a document.
It ends a line, not a paragraph:
the next line has no indent and no gap before it,
so the text looks like a paragraph but is not one,
and \hologo{LaTeX} cannot fill the broken line evenly any more.
At the end of a paragraph, or twice in a row to fake an empty line,
it produces the warning
\begin{codeblock}[style=shell]
Underfull \hbox (badness 10000) in paragraph at lines 12--14
\end{codeblock}
and a paragraph whose last line is stretched across the full width.
The habit comes from the Enter key of a word processor; in \hologo{LaTeX}, a new paragraph is an empty line.

\code|\\| does belong in a few places:
the rows of a table (\autoref{sec:tables}), the lines of a multi-line equation (\autoref{sec:align}),
and the title or author block, to break a long title by hand.

\subsection{More space, or a new page}

Space between paragraphs, or a new page, is mostly left to the document class.
When you do need it:

\begin{itemize}
    \item \code|\smallskip|, \code|\medskip| and \code|\bigskip| add vertical space between paragraphs;
        \code|\vspace{1cm}| adds a fixed amount.
    \item \code|\newpage| starts a new page.
        \code|\clearpage| does the same, but first places all figures and tables that are still waiting (\autoref{sec:floats}).
    \item \code|\quad| and \code|\qquad| are wide horizontal spaces;
        \code|\,| is a thin one, e.g.\ between a number and its unit, \code|250\,ml| (\autoref{sec:siunitx} has a better way for units).
\end{itemize}

\subsection{Space after a full stop}

In English, \hologo{LaTeX} puts slightly more space after a full stop that ends a sentence than between words.
It guesses: a full stop after a lower-case letter ends a sentence, a full stop after a capital letter does not.
Where the guess is wrong, help it:

\begin{itemize}
    \item After an abbreviation in the middle of a sentence, a backslash and a space give a normal space:
        \code|e.g.\ a mug|.
        If the abbreviation belongs to the next word, use the tilde: \code|Dr.~Watson|.
    \item A sentence that ends with a capital letter needs \code|\@| before the full stop:
        \code|measured by NASA\@.|
\end{itemize}

The difference is small, but it is visible in a long text.
With \code|\frenchspacing| in the preamble, all spaces are the same and the question disappears;
that is the usual setting for German, and \texttt{babel} switches it on for German texts (\autoref{sec:babel}).

\section{Quotes, dashes and special characters}\label{sec:special-chars}

\subsection{Quotation marks}

\hologo{LaTeX} makes opening and closing quotation marks from two different characters:
the backtick \code|`| (on most keyboards next to the~1 or the Backspace key) for the opening one,
and the apostrophe \code|'| for the closing one.
Doubled, they give double quotation marks.
The straight double quote \code|"| gives the closing mark at both ends:

\owncode{ch04-quotes.tex}
\ownrender{ch04-quotes.pdf}

The package \texttt{csquotes} does this for you with \code|\enquote|,
and chooses the marks of the document's language
(with \texttt{babel} set to German: \quotedblbase like this\textquotedblleft):

\owncodetitled{ch04-csquotes-preamble.tex}{Preamble}
\owncodetitled{ch04-csquotes.tex}{Document}
\ownrender{ch04-csquotes.pdf}

Typographic quotation marks pasted from another program (\texttt{“...”}) also work in the text,
but not inside commands that expect plain characters, e.g.\ in a file name or a label.

\subsection{Dashes}

There are three kinds of dash, typed as one, two or three hyphens:

\begin{itemize}
    \item \code|-|: the hyphen, inside words: \enquote{well-known}.
    \item \code|--|: the en dash, for ranges: \enquote{pages 10--20}.
        In British typesetting, also as a dash in a sentence, with spaces around it: \enquote{a pause -- like this}.
    \item \code|---|: the em dash, a dash in a sentence in American typesetting, without spaces.
\end{itemize}

\owncode{ch04-dashes.tex}
\ownrender{ch04-dashes.pdf}

A minus sign is none of them: it belongs in math mode, \code|$-5$| (\autoref{chap:math}).
Dashes pasted from a word processor (\texttt{–}, \texttt{—}) work in the text, like the quotation marks,
but they break file names, labels and options that expect a plain hyphen (\autoref{sec:common-errors}).

\subsection{Other characters}

Since 2018, \hologo{LaTeX} reads source files as UTF-8.
Letters with accents and umlauts can be typed directly: \enquote{Zürich}, \enquote{café}, \enquote{Straße}.
Together with \code|\usepackage[T1]{fontenc}| in the preamble (\autoref{sec:preamble}),
such words are also hyphenated correctly and can be searched and copied from the PDF.

A few characters have a command of their own, e.g.\
\code|\ldots| for the ellipsis (\ldots),
\code|\S| for the section sign (\S),
\code|\texteuro| for the euro sign (\texteuro)
and \code|\copyright| (\copyright).
Characters that \hologo{LaTeX} does not know at all,
e.g.\ an arrow ($\rightarrow$) or a \enquote{less than or equal} ($\leq$) pasted into the text,
stop the compilation with \texttt{Unicode character ... not set up for use with LaTeX}.
Most of them are math symbols and belong in math mode (\autoref{sec:symbols}).

\section{Lists}\label{sec:lists}

\hologo{LaTeX} has three kinds of lists, all with the same shape:
\code|\begin|, one \code|\item| per entry, \code|\end|.

\begin{tcbraster}[raster columns=3, raster equal height, raster column skip=6pt,
        raster left skip=-\exoverhang, raster right skip=-\exoverhang]
    \tcbinputlisting{codebox, grow to left by=0pt, grow to right by=0pt, title=itemize,
        listing file={\slidesdir/listings/02-itemize.tex}, listing options={style=code, numbers=none}}
    \tcbinputlisting{codebox, grow to left by=0pt, grow to right by=0pt, title=enumerate,
        listing file={\slidesdir/listings/02-enumerate.tex}, listing options={style=code, numbers=none}}
    \tcbinputlisting{codebox, grow to left by=0pt, grow to right by=0pt, title=description,
        listing file={\slidesdir/listings/02-description.tex}, listing options={style=code, numbers=none}}
\end{tcbraster}
\sliderender{02-lists.pdf}

\begin{description}
    \item[\texttt{itemize}] A list with bullets, where the order does not matter.
    \item[\texttt{enumerate}] A numbered list, where the order matters.
    \item[\texttt{description}] Terms and their definitions.
        The term goes in the optional argument of \code|\item|.
\end{description}

Lists can be nested, up to four levels deep.
The labels change with the level by themselves:

\owncode{ch04-nested.tex}
\ownrender{ch04-nested.pdf}

The optional argument of \code|\item| replaces the label of a single entry:

\owncode{ch04-item-label.tex}
\ownrender{ch04-item-label.pdf}

To change the labels or the spacing of all lists, use the package \texttt{enumitem} (\autoref{sec:enumitem}).
An empty line inside a list does no harm; it only starts a new paragraph within the current item.

\section{The preamble}\label{sec:preamble}

The preamble is the part of the source file between \code|\documentclass| and \code|\begin{document}|.
It holds the settings for the whole document;
nothing in it is printed.
A typical preamble loads a few packages:

\slidecode{02-preamble.tex}

\begin{description}
    \item[\texttt{fontenc}] With the option \texttt{T1}, the fonts contain letters with accents and umlauts as one character,
        so such words are hyphenated correctly and can be searched and copied from the PDF.
    \item[\texttt{babel}] The language of the document: hyphenation rules and words such as \enquote{Contents} or \enquote{Figure}
        (\autoref{sec:babel}).
    \item[\texttt{parskip}] A gap between paragraphs instead of an indent (\autoref{sec:parskip}).
    \item[\texttt{amsmath}] Everything for math (\autoref{chap:math}).
\end{description}

Whenever you wonder \enquote{can \hologo{LaTeX} do this?}, the answer is almost always yes, with one line in the preamble.
Over time, your preamble grows into a collection of such lines,
and the next report starts with a copy of it.
\autoref{sec:preamble-file} shows how to keep it in a file of its own.

\section{Packages}\label{sec:packages}

A package adds commands, or changes how the document looks.
It is loaded in the preamble with \code|\usepackage|,
which has the same shape as \code|\documentclass|:

\begin{codeblock}
\usepackage[english]{babel}
\end{codeblock}

The options go in the square brackets, the name of the package in the braces.
Several packages without options can share one command,
e.g.\ \code|\usepackage{booktabs, graphicx}|,
but one package per line is easier to read and to comment.

\begin{itemize}
    \item \hologo{TeX} Live, Mac\hologo{TeX} and Overleaf come with almost all packages.
        \hologo{MiKTeX} installs a package the first time a document loads it (\autoref{sec:miktex}).
    \item There are about 6\,000 packages on \gls{ctan}.
        Every one of them has a page at \texttt{ctan.org/pkg/<name>}, e.g.\ \myhref{https://ctan.org/pkg/booktabs}{ctan.org/pkg/booktabs},
        with its documentation.
        On your own installation, \code|texdoc booktabs| in a terminal opens the same documentation.
    \item The order of the packages rarely matters.
        The important exception is \texttt{hyperref}, which goes last (\autoref{sec:hyperref}).
    \item A package that is not installed stops the compilation with
        \texttt{File `name.sty' not found} (\autoref{sec:common-errors}).
\end{itemize}

\autoref{chap:packages-chapter} describes the packages most reports need.

\section{Paragraph spacing}\label{sec:parskip}

By default, every paragraph after the first one in a section is indented, and there is no gap between paragraphs:

\slidecodetitled{02-paragraphs.tex}{Document}
\sliderender{02-paragraphs.pdf}

That is fine for a novel, but unusual for a report.
The fix most beginners find is \code|\noindent| at the start of every paragraph:

\slidecodetitled{02-noindent.tex}{Document}
\sliderender{02-noindent.pdf}

It removes the indent, but now the paragraphs run into each other,
and every paragraph where you forget it looks broken.
The better fix is one line in the preamble, the package \texttt{parskip}:

\slidecodetitled{02-parskip.tex}{Preamble}
\slidecodetitled{02-paragraphs.tex}{Document}
\sliderender{02-paragraphs-parskip.pdf}

No indent, a gap between paragraphs, and the document itself is back to plain paragraphs.
To undo it, delete one line.
\texttt{parskip} also adjusts the spacing in lists and in the table of contents,
which setting \code|\parindent| and \code|\parskip| by hand would not.

Gaps between paragraphs are the usual convention in German-language documents,
so for much of your coursework \texttt{parskip} is simply the right choice.
In general: if you find yourself repeating a command everywhere, look for a package or a setting that does it once.

\section{Setting the look of emphasis}\label{sec:emph-preamble}

\code|\emph| (\autoref{sec:emph}) marks text as emphasised and leaves the look to the document.
To make all emphasised text bold instead of italic, redefine \code|\emph| in the preamble:

\slidecodetitled{02-emph-preamble.tex}{Preamble}
\slidecodetitled{01-emph.tex}{Document}
\sliderender{02-emph.pdf}

The document did not change, only the line in the preamble,
and every \code|\emph| in the document follows.
\code|\renewcommand| replaces an existing command,
and \code|#1| stands for its argument
(\autoref{sec:own-commands} explains both).
Bold is only an example; what matters is that the decision lives in one place.
